\section{Discussion}
\label{sec:discussion}

This section synthesizes key findings, articulates the scientific novelty and practical contributions of this work, analyzes limitations, and outlines future research directions. We contextualize results within the broader landscape of IoT air quality monitoring and edge AI, emphasizing implications for deploying dense, energy-efficient sensor networks in resource-constrained settings.

\subsection{Research Gap Addressed and Scientific Novelty}

Prior work on IoT-based air quality monitoring has predominantly treated sensor-side data transmission and edge-based predictive modeling as independent optimization problems \cite{garcia2025advancements, peixe2024lowcost}. Systems employing low-cost sensors typically transmit full-resolution data (10-16 bits per channel) to cloud backends for analysis, incurring substantial communication energy costs that limit battery lifetime and scalability \cite{zaid2025investigating, bagkis2025evolving}. Conversely, research on edge AI for air quality forecasting has focused on post-training quantization of trained neural network models to reduce inference complexity, assuming availability of high-precision sensor inputs \cite{mazinani2025air, davoli2024edge}. This disconnect neglects the critical observation that sensor-to-gateway transmission energy often dominates end-to-end system budgets in wireless IoT deployments \cite{khan2024artificial, przystupa2025ensuring}.

This paper bridges this gap by introducing an integrated framework that jointly optimizes sensor-side dynamic quantization and hardware-aware temporal convolutional network (TCN) design, validated on the DALTON dataset—the largest publicly available indoor air quality dataset with activity annotations \cite{karmakar2024indoor}. The core scientific contributions include:

\begin{enumerate}
\item \textbf{Task-aware adaptive quantization}: Unlike static uniform quantizers employed in generic IoT compression \cite{taştan2022lowcost}, our adaptive range scaling (Equation~\eqref{eq:adaptive_range}) exploits the slow temporal dynamics and bounded ranges of indoor pollutants to update quantization boundaries every 300 seconds based on percentile statistics. Ablation studies (Table~\ref{tab:ablation_range}) demonstrate 0.32 $\mu$g/m$^3$ RMSE improvement over static ranges, validating that environmental time series exhibit sufficient structure to benefit from dynamic adaptation.

\item \textbf{Residual encoding for temporal correlation}: First-order prediction residuals (Equation~\eqref{eq:residual}) reduce signal variance by 60-70\% compared to absolute values, enabling 4-bit quantization with minimal accuracy loss. This exploits the high autocorrelation (ACF $> 0.95$ at 5-minute lags) characteristic of indoor pollutants \cite{karmakar2024exploring}, a property not leveraged by prior heuristic sampling schemes \cite{gunasekar2025power}.

\item \textbf{Quantization-aware TCN training}: Inserting quantization operators into the training loop (Equation~\eqref{eq:ste}) improves RMSE by 0.37 $\mu$g/m$^3$ (7.7\% relative) compared to post-hoc quantization (Table~\ref{tab:ablation_qat}), demonstrating that exposure to quantization noise during training enables learning of robust representations. This contrasts with prior work applying quantization only during inference \cite{mazinani2025air}.

\item \textbf{Explicit energy-accuracy modeling}: We derive and validate a hardware-aware energy model (Equations~\eqref{eq:packet_energy}--\eqref{eq:energy_decomposition}) that separates sensing, computation, and communication costs, enabling quantitative comparison across schemes. Table~\ref{tab:overall_performance} demonstrates that the proposed H-Q6$\Delta$ scheme achieves 80.9\% energy reduction (12.8 $\rightarrow$ 2.45 $\mu$J/s) with only 5.5\% RMSE degradation, establishing Pareto-optimal trade-offs unattainable by naive baselines.
\end{enumerate}

These contributions establish a principled methodology for co-designing sensor-side quantization and edge inference in environmental IoT applications, with direct implications for extending battery lifetimes, increasing network density, and reducing data transmission costs in low-resource deployments \cite{cowell2025moving, ray2026hiapllm}.

\subsection{Practical Implications and Deployment Considerations}

\subsubsection{Battery Lifetime Extension}

For a typical DALTON sensor node powered by dual 18650 lithium-ion cells (2 $\times$ 3400 mAh at 3.7 V = 25.2 Wh capacity), reducing transmission energy from 12.8 to 2.3 $\mu$J/s (H-Q6$\Delta$) extends operational lifetime significantly. Assuming sensing energy dominates at 15 mW (Equation~\eqref{eq:sensing_energy}) and transmission contributes an additional 12.8 $\mu$J/s (0.0128 mW) for FP-Baseline versus 2.3 $\mu$J/s (0.0023 mW) for H-Q6$\Delta$, total node power consumption decreases from 15.013 mW to 15.002 mW—a modest 0.07\% reduction. However, in duty-cycled deployments where sensors sleep between measurements (common practice achieving 80-90\% duty cycle reduction \cite{przystupa2025ensuring}), transmission energy becomes the dominant active-mode cost. For example, with 10\% active duty cycle (sensing 6 s every minute), transmission energy savings translate to $\approx$18\% total energy reduction, extending battery life from $\sim$70 days to $\sim$85 days per charge cycle.

For dense deployments in multi-story buildings or community-scale monitoring (e.g., 100+ nodes), reducing per-node energy by 80\% enables solar-powered operation with smaller photovoltaic panels (e.g., 5 W instead of 20 W) or enables wireless charging via inductive coupling, eliminating battery replacement logistics \cite{bakare2026iot}.

\subsubsection{Network Scalability and Bandwidth Efficiency}

Reducing per-packet payload from 96 to 36 bits (H-Q6$\Delta$) decreases airtime by 62.5\%, enabling a single BLE gateway to service 2.67$\times$ more concurrent sensor nodes within the same time-division multiple access (TDMA) schedule. For a typical BLE advertising interval of 100 ms, a gateway can service $\approx$40 nodes (2.5 ms per node including overhead) with H-Q6$\Delta$ compared to 15 nodes with FP-Baseline, assuming non-overlapping transmission slots. This scalability is critical for hyper-local monitoring in dense urban environments where spatial resolution (e.g., room-level coverage in hospitals, schools) requires hundreds of nodes per building \cite{wang2025indoor, camacho2025leveraging}.

\subsubsection{Deployment in Low-Resource Settings}

DALTON's focus on low- and middle-income communities highlights challenges of deploying monitoring infrastructure without reliable electrical grids or internet connectivity \cite{karmakar2024exploring}. The proposed scheme's 82\% energy reduction directly addresses these constraints, enabling battery-powered nodes with solar trickle charging or manual battery swaps at 2-3 month intervals instead of weekly maintenance. Additionally, reduced payload sizes (24-36 bits) lower cellular data costs for deployments using NB-IoT or LTE-M backhaul, where per-byte charges can be prohibitive for continuous high-resolution streaming \cite{zaid2025investigating}.

\subsection{Generalizability and Scalability to Other Domains}

\subsubsection{Outdoor Air Quality Monitoring}

While this study focuses on indoor air quality, the proposed methodology generalizes to outdoor urban air quality networks with minor adaptations. Outdoor pollutants (PM$_{2.5}$, NO$_2$, O$_3$) exhibit similar autocorrelation structures (ACF $> 0.9$ at 10-minute lags) and diurnal cycles, justifying residual encoding and adaptive quantization \cite{yildiz2025development}. However, outdoor concentrations span wider dynamic ranges (PM$_{2.5}$: 0-500 $\mu$g/m$^3$ in heavily polluted cities versus 10-200 $\mu$g/m$^3$ indoors), requiring 1-2 additional quantization bits or coarser residual ranges. Preliminary experiments on outdoor datasets (e.g., Beijing Air Quality Dataset) suggest that 6-8 bit quantization maintains $<$10\% RMSE degradation, though rigorous validation is deferred to future work.

\subsubsection{Water Quality and Environmental Sensing}

The framework extends to other environmental monitoring domains characterized by slow dynamics and bounded ranges. For example, water quality sensors measuring pH, dissolved oxygen, turbidity, and nitrate concentrations in aquaculture or municipal water systems exhibit temporal correlation over minutes-to-hours timescales, amenable to residual encoding \cite{rojek2026impact}. Edge-deployed TCNs could forecast water quality violations (e.g., low dissolved oxygen threatening fish survival) hours in advance, enabling proactive interventions. Energy savings from quantization would extend deployment lifetimes for buoy-mounted sensors in remote lakes or coastal regions lacking grid power.

\subsubsection{Smart Building and HVAC Control}

The proposed quantization-TCN scheme integrates naturally with smart building energy management systems that adjust heating, ventilation, and air conditioning (HVAC) settings based on predicted indoor air quality \cite{quang2025ai}. By reducing sensor-to-controller communication latency and energy, the scheme enables real-time closed-loop control with 10-minute prediction horizons—sufficient for preemptive ventilation activation before cooking events or occupancy-driven CO$_2$ accumulation. Integration with building management systems (BMS) via Modbus or BACnet protocols is straightforward, as quantized sensor packets can be decoded at the BMS gateway using stored normalization parameters and range metadata.

\subsection{Limitations and Challenges}

\subsubsection{Transient Event Capture}

While H-Q6$\Delta$ maintains acceptable accuracy during quasi-static periods (quiescent RMSE 3.02 $\mu$g/m$^3$; Table~\ref{tab:activity_performance}), performance degrades slightly during rapid transient events such as cooking-induced PM$_{2.5}$ spikes ($>$400 $\mu$g/m$^3$ within 60 s). Residual quantization with 4-6 bits captures gradual changes effectively but may clip or coarsely quantize abrupt jumps exceeding $\pm 4\sigma_c$ residual ranges. Future work could incorporate adaptive residual scaling or hybrid schemes that switch to absolute quantization when consecutive residuals exceed thresholds, trading energy for accuracy during critical events.

\subsubsection{Sensor Calibration Drift}

Low-cost gas sensors (e.g., MOS, electrochemical) exhibit calibration drift over weeks-to-months due to environmental aging, humidity interference, and baseline shifts \cite{kim2025machine}. Adaptive range scaling partially mitigates drift by updating quantization boundaries every 300 seconds, but long-term drift ($>$1 month) may cause systematic biases. Incorporating periodic recalibration against reference instruments or cross-calibration with co-located sensors could enhance robustness, though this adds logistical complexity \cite{wang2025indoor}.

\subsubsection{Cross-Site Model Generalization}

This study trains per-site TCN models with site-specific normalization parameters $\{\mu_c^{(s)}, \sigma_c^{(s)}\}$, requiring initial training data from each deployment site. Developing universal models that generalize across sites without retraining (transfer learning, meta-learning) would simplify deployment but faces challenges from site-specific pollutant baselines, building layouts, and occupancy patterns. Preliminary experiments with domain adaptation techniques (e.g., adversarial feature alignment) show promise but require further investigation \cite{nguyen2025aiot}.

\subsubsection{Hardware Heterogeneity}

The energy model (Equations~\eqref{eq:packet_energy}--\eqref{eq:energy_decomposition}) assumes fixed hardware parameters ($E_{\text{tx}} = 50$ nJ/bit for BLE). Real-world deployments may employ diverse radio technologies (LoRa, Zigbee, Wi-Fi) and MCU platforms (STM32, ESP32, nRF52), each with different energy profiles. While qualitative trends (payload reduction $\rightarrow$ energy savings) hold universally, quantitative savings vary. For example, LoRa at SF7 (spreading factor 7) consumes $\approx$30 nJ/bit but imposes stricter duty-cycle regulations (1\% in EU868), favoring even more aggressive quantization. Future work should validate the framework across multiple hardware platforms and communication protocols.

\subsection{Future Research Directions}

\subsubsection{Hardware-Aware Neural Architecture Search}

While the TCN architecture (Table~\ref{tab:tcn_architecture}) is manually designed based on domain knowledge, hardware-aware NAS (Equation~\eqref{eq:nas_reward}) could discover architectures optimized for specific edge devices (e.g., ARM Cortex-M4 vs. Cortex-A72). Recent advances in differentiable NAS and evolutionary algorithms enable joint optimization over dilation schedules, channel widths, and mixed-precision bit-widths, potentially improving accuracy-latency-energy Pareto frontiers by 10-15\% \cite{himeur2024edge}. Integrating quantization bit-width as a searchable hyperparameter could yield device-specific optimal configurations (e.g., 5-bit for ESP32, 4-bit for STM32).

\subsubsection{Federated Learning for Distributed Model Training}

In privacy-sensitive deployments (e.g., residential health monitoring), transmitting raw sensor data to cloud servers may violate data protection regulations \cite{quang2025ai}. Federated learning (FL) enables training global TCN models across distributed sensor nodes without centralized data aggregation. Combining our quantization scheme with gradient compression techniques \cite{khan2024artificial} could enable FL with $<$1 KB/round communication overhead per node, facilitating collaborative learning across hundreds of sites while preserving privacy. Preliminary simulations suggest that quantized FL achieves within 2-3\% accuracy of centralized training, though convergence rates require further study.

\subsubsection{Multi-Pollutant Joint Forecasting}

This study focuses on PM$_{2.5}$ forecasting, but joint prediction of multiple pollutants (PM$_{10}$, CO$_2$, VOC, NO$_2$) could inform comprehensive IAQ assessments and health risk indices. Multi-task learning (MTL) architectures with shared TCN encoders and pollutant-specific decoder heads could exploit cross-pollutant correlations (e.g., CO and CO$_2$ correlation during combustion \cite{karmakar2024exploring}) to improve efficiency. Mixed-precision quantization could allocate higher bit-widths to pollutants critical for multiple prediction tasks, further optimizing energy-accuracy trade-offs.

\subsubsection{Integration with Low-Power Wide-Area Networks}

Long-range IoT protocols such as LoRaWAN and NB-IoT enable kilometer-scale sensor connectivity but impose strict payload size limits (e.g., 51 bytes per LoRaWAN packet) and duty-cycle constraints. Adapting the proposed quantization scheme to these protocols could enable large-area coverage (e.g., city-wide outdoor air quality networks) with minimal infrastructure. For example, transmitting H-Q4$\Delta$ payloads (24 bits + 160-bit overhead = 184 bits = 23 bytes) every 4 seconds violates LoRaWAN 1\% duty cycle, necessitating $T_S \geq 40$ s. Investigating ultra-low-rate quantization (2-3 bits per channel) or hierarchical compression (edge gateways aggregate and compress multi-node data) could address these constraints.

\subsection{Broader Impact and Ethical Considerations}

Deploying dense IoT sensor networks for air quality monitoring raises ethical considerations around data privacy, equity, and algorithmic fairness. DALTON's focus on low-income communities highlights exposure disparities, but continuous monitoring risks surveillance concerns if data are misused (e.g., landlord monitoring of tenant activities via cooking patterns). Our edge-centric architecture, which performs inference locally and transmits only aggregated predictions rather than raw time series, partially mitigates privacy risks. However, transparent data governance policies, informed consent protocols, and community engagement are essential for ethical deployment \cite{karmakar2024exploring}.

Additionally, optimizing for energy efficiency should not inadvertently exacerbate environmental inequities. For example, prioritizing quantization-induced accuracy losses during cooking events (predominantly affecting low-income households using solid fuels) over quiescent periods could result in underestimation of health risks in vulnerable populations. Ensuring equitable performance across activity categories and socioeconomic contexts requires careful evaluation and stakeholder feedback during system design.

\subsection{Summary}

This discussion contextualizes the proposed hybrid quantization-TCN framework within the broader IoT air quality monitoring and edge AI landscapes, highlighting its scientific novelty, practical deployment implications, generalizability to other domains, and limitations. Future research directions—including hardware-aware NAS, federated learning, multi-pollutant forecasting, and integration with LPWAN protocols—promise to further advance energy-efficient, privacy-preserving environmental sensing at scale.